% 复刻学校的中期报告样例
\section{课题主要研究内容及进度}
\subsection{课题主要研究内容}
本课题将从以下几个方面开展研究：

1. 人工心脏泵体数学模型的建立。……

2. 人工心脏泵体的优化设计。……

3. ……

\subsection{进度介绍}
目前，已经完成了人工心脏泵体数学模型的建立，人工心脏泵体的优化设计。正在运用 CFD 软件进行相关结果的验证。

……

\section{已完成的研究工作及结果}
\subsection{人工心脏泵体的设计}
\subsubsection{叶片的参数设计}
……

结合计算的结果和结构设计的需要，调整了部分参数，基本参数如表~\ref{tab:blade-parameters} 所示。

\begin{table}[htbp]
  \centering
  \caption{叶片主要参数表}
  \label{tab:blade-parameters}
  \begin{tabularx}{\linewidth}{*{3}{>{\centering\arraybackslash}X}}
    \toprule
    参数 & 数值 & 单位 \\
    \midrule
    比转速 ns & 71 & —— \\
    叶轮进口直径 & 15 & mm \\
    叶轮出口直径 & 38 & mm \\
    叶轮进口厚度 & 4 & mm \\
    叶轮出口厚度 & 4 & mm \\
    叶轮进口安装角 & 23.6 & 度 \\
    叶轮出口安装角 & 29.0 & 度 \\
    包角 & 100 & 度 \\
    \bottomrule
  \end{tabularx}
\end{table}

\subsubsection{叶片的三维造型}

\subsubsection{涡室的参数设计}
不加导叶的压水室大体上可分为螺旋形压水室和环形压水室两种。……

插图格式如图~\ref{fig:example-image} 所示。

\begin{figure}[htbp]
  \centering
  % 使用模板自带的示例图片，撰写时替换为自己的图片。
  \includegraphics[width=.3\linewidth]{golfer}
  \caption{示例图片}
  \label{fig:example-image}
\end{figure}
\FloatBarrier

隔舌一般指的是涡室螺旋部分的始端。第八断面与隔舌之间的夹角被称为隔舌安放角。按照经验选取，本文 $\alpha_0=20^\circ$。

圆形截面的计算公式如下，
\begin{gather}
  \mathit{major}=\mathit{minor}+\frac{b_3}{2}
    +\frac{\left|\frac{b_3}{2}-\mathit{minor}\right|}{2}
  \label{eq:volute-major}
  \\
  \mathit{height}=\mathit{minor}
    +\sqrt{(\mathit{major})^2-\left(\frac{b_3}{2}\right)^2}
  \label{eq:volute-height}
\end{gather}

其中，major，minor 对于椭圆段分别指的是长短轴；对于圆形段指的是半径。

……

\subsection{人工心脏叶轮性能优化算法}
\subsubsection{以降低血液中的切应力为目标建立优化模型}
按照血液在运动过程中，由体积相等，可得 5L 液体的作用时间的表达式如下，
\begin{equation}
  t=\frac{\pi b}{4Q}\left(D_2^2-D_1^2\right)
  \label{eq:blood-transit-time}
\end{equation}
\noindent 式中\quad
\begin{tabular}[t]{@{}l@{}}
  $t$——作用时间，s；\\
  $Q$——流量，$\mathrm{m}^3/\mathrm{s}$；\\
  $b$——叶片的宽度，m。
\end{tabular}

……

\section{后期拟完成的研究工作及进度安排}
\subsection{后期拟完成的研究工作}
后续需要完成的工作主要是……。

\subsection{进度安排}
具体进度安排如下：

2022 年 4 月 1 日——2022 年 4 月 15 日：整理已经做完的数据，分析……。

2022 年 4 月 16 日——2022 年 4 月 26 日：整理完成……。

2022 年 4 月 27 日——2022 年 5 月 15 日：……内容的研究，并完成……。

2022 年 5 月 16 日——2022 年 6 月 20 日：……。

\section{存在的困难及解决方案}
\subsection{存在的困难}
1. 仿真前需要进行……。

2. ……

3. 有些实验设备存在购买滞后……。

……

\subsection{解决方案}
针对目前存在的问题和困难，提出以下解决方案：

1. ……

2. ……

3. ……

……

\section{论文按时完成的可能性}
本课题目前已经完成了……。

目前正在进行泵体内不同工况下切应力分布图的仿真结果的统计，将要进行……。

综上所述，论文能够按期完成，并取得一定的研究成果。

% Local Variables:
% TeX-master: "../report"
% TeX-engine: xetex
% End:
